%% kapa: exam=1 round=28 year=2017 period=2 form=A subject=appraisal_law subject_ko=감정평가관계법규 seq=23 number=23
%% source: Q-Net 감정평가사 기출문제 (cst00309) artlSeq=5199603 file=extracted/1008/5199603/2017년 제28회 감정평가사 1차시험 문제지/2교시 A형.pdf
%% conversion: kordoc 4.24.51 -> tools/gen1.py (auto); review=proofread-sonnet; flags=none
\begin{question}{23}
국유재산법령상 행정재산의 사용허가의 취소ㆍ철회사유에 해당하지 않는 것은?
\begin{choices}
  \item 해당 재산의 보존을 게을리한 경우
  \item 부실한 증명서류를 제시하여 사용허가를 받은 경우
  \item 중앙관서의 장이 사용허가 외의 방법으로 해당 재산을 관리ㆍ처분할 필요가 있다고 인정되는 경우
  \item 납부기한까지 사용료를 납부하지 않은 경우
  \item 중앙관서의 장의 승인 없이 사용허가를 받은 재산의 원래 상태를 변경한 경우
\end{choices}
\end{question}
\answer{3}
